% ch03-requirements.tex – Chapter 3: requirements specification (final round, Writer).
\ChapterMission{המשימה: לתרגם את יעדי התכן לדרישות שלכל אחת יש יעד מספרי ושיטת אימות}
\ChapterIcon{signal}
\chapter{מפרט הדרישות}\label{chap:requirements}
\begin{scenebox}
ערב. פועל בסרבל כתום מטפס על סולם ליד עמוד הרמזור ומכוון קופסה קטנה אל המדרכה.
מהנדסת עומדת מתחתיו עם מחשב נייד ובוחנת קו על המסך. אדם עובר באזור ההמתנה, והקו קופץ.
היא רושמת משהו ומבקשת ממנו לחזור על המעבר.
\end{scenebox}

הכיול שבתמונה הוא בדיוק מה שהפרק דורש: כל דרישה נבדקת, ודרישה בלי שיטת אימות היא משאלה בלבד.
הפרק מתרגם את היעדים של פרק~\ref{chap:problem} לדרישות פונקציונליות, לדרישות שאינן פונקציונליות ולאילוצים, ומגדיר את הגבולות הבטוחים שבתוכם מותר לבקר לקצר את האדום.

\section{מקורות הדרישות}
הדרישות נשענות על שלושה מקורות.
לחישה בשפת המדרכה קיים מפרט רכש עירוני מפורט, שמגדיר אזור המתנה, זמן זיהוי, עבודה בתאורה חלשה והתנהגות בתקלה \cite{tfl2015kerbside}.
לתזמון האות, התקן האמריקני קובע את משך מופע ההליכה (\en{WALK}), את מהירות ההליכה שלפיה מחושב זמן הפינוי ואת מרווח הביטחון שאחריו, ומתיר ספירה לאחור רק במופע הפינוי \cite{fhwa2025mutcd}.
לירוק המזערי לכלי רכב, מדריך התזמון של מינהל הדרכים הפדרלי בארצות הברית (\en{FHWA}) מציג טווחים מקובלים לפי סוג הדרך \cite{fhwa2008stm}.
שני האחרונים הם מסמכים אמריקניים; התאמתם לתקן הישראלי היא חלק מחבילת העבודה הראשונה (פרק~\ref{chap:workplan}).

\section{הגבולות הבטוחים של האדום המסתגל}
האדום המסתגל מקצר רק את המתנת הולך הרגל, כלומר את הירוק של כלי הרכב החוצים. הוא אינו נוגע במופע ההליכה, בפינוי או במרווח הביטחון.

\begin{table}[!ht]
  \centering
  \caption{מפרט הדרישות. סיווג: פ = פונקציונלית, ל = לא פונקציונלית, א = אילוץ. יעדים ללא מקור הם יעדי הצוות}
  \label{tab:req}
  \small
  \begin{tabularx}{\linewidth}{lXcXX}
    \toprule
    \headrow \thead{מזהה} & \thead{דרישה} & \thead{סוג} & \thead{יעד} & \thead{אימות} \\
    \midrule
    \en{REQ-01} & זיהוי הולך ממתין באזור ההמתנה, עומד או נע & פ & תוך \unit{500}{ms}; אזור ברוחב \unit{2.4\textendash 10}{m} \cite{tfl2015kerbside} & בדיקת שטח מול וידאו מתויג \\
    \en{REQ-02} & ספירת הממתינים ומדידת זמן ההמתנה & פ & ספירה נכונה ב-\pct{90} מן המופעים; שגיאת זמן עד \unit{2}{s} & השוואה לתיוג ידני של \num{500} מופעים \\
    \en{REQ-03} & פעולה בלילה ובתאורה חלשה, דחיית צללים וכלי רכב & ל & עד \unit{10}{lux} \cite{tfl2015kerbside} & בדיקת לילה בשטח \\
    \en{REQ-04} & עדיפות להולך יחיד: בקשה מוקדמת למופע הולכי רגל & פ & בקשה תוך \unit{1}{s} מחציית סף ההמתנה & סימולציה ויומן בקר \\
    \en{REQ-05} & ירוק מזערי לכלי רכב אינו מקוצר & א & לפי סוג הדרך \cite{fhwa2008stm}; אפס הפרות & ביקורת יומני בקר \\
    \en{REQ-06} & מופע הליכה אינו מקוצר & א & \unit{7}{s} לפחות \cite{fhwa2025mutcd} & בדיקת קבלה \\
    \en{REQ-07} & פינוי ומרווח ביטחון אינם מקוצרים & א & משוואה~\eqref{eq:clearance}; מרווח של \unit{2}{s} לפחות \cite{fhwa2025mutcd} & בדיקת קבלה \\
    \en{REQ-08} & קיצור מרבי ושמירה על המחזור והתיאום & א & עד \pct{30} מן האדום הנומינלי; עיכוב רכב נוסף עד \pct{5} & סימולציה ומדידה בפיילוט \\
    \en{REQ-09} & ספירה לאחור רק במופע הפינוי, חשוכה אחרי האפס & פ & בכל מופע פינוי \cite{fhwa2025mutcd} & בדיקת כל המופעים בקבלה \\
    \en{REQ-10} & מצב כשל בטוח: חזרה לתזמון קבוע & ל & תוך \unit{3}{s} מתקלה \cite{tfl2015kerbside} & הזרקת תקלות \\
    \en{REQ-11} & פרטיות: עיבוד בקצה, בלי שמירת תמונות & א & אפס תמונות שמורות או משודרות & ביקורת קוד ואחסון \\
    \en{REQ-12} & זמינות המערכת & ל & \pct{99} משעות הפעולה & יומני ניטור בפיילוט \\
    \en{REQ-13} & קמפיין: תצוגה באתר ומסרים בתקשורת & פ & \pct{60} מן הנשאלים באזור מכירים את המסר & סקר מודעות \\
    \bottomrule
  \end{tabularx}
\end{table}
זמן הפינוי נגזר מאורך המעבר ומהירות ההליכה (משוואה~\eqnum{eq:clearance}):
\begin{equation}
  t_{\hebmath{פינוי}} = \frac{L}{v}, \qquad v \le 1.07~\mathrm{m/s}
  \label{eq:clearance}
\end{equation}
כאן $L$ הוא אורך המעבר עד המדרכה שממול או עד אי תנועה, ו-$v$ היא מהירות ההליכה. התקן קובע \unit{3.5}{ft/s}, שהם כ-\unit{1.07}{m/s}, ומחייב מהירות נמוכה יותר כשהולכים איטיים או משתמשי כיסא גלגלים חוצים במקום באופן קבוע \cite{fhwa2025mutcd}.
אחרי הפינוי בא מרווח ביטחון של \unit{2}{s} לפחות, ומופע ההליכה עצמו נמשך \unit{7}{s} לפחות, או \unit{4}{s} כשהתנאים מתירים \cite{fhwa2025mutcd}.
בצד כלי הרכב, ירוק מזערי נמוך מדי מגדיל את הסיכון לתאונות מאחור; בעורק ראשי מקובל ירוק מזערי של \unit{7\textendash 15}{s} לתנועה הישרה \cite{fhwa2008stm}.
טבלה~\ref{tab:req} מרכזת את כל הדרישות. הדרישות \en{REQ-05} עד \en{REQ-08} הן הגבולות הבטוחים, ואף אחת מהן אינה ניתנת להתפשרות לטובת קיצור ההמתנה.

\section{הערות לדרישות}
דרישת הספירה (\en{REQ-09}) חשובה במיוחד משום שהתקן מזהיר שלעיתים הספירה מחושבת ביחידה נפרדת מן הבקר \cite{fhwa2025mutcd}; כשהאדום המסתגל משנה את אורכי המופעים, ראש הספירה חייב לקבל את התזמון מן הבקר, ולא לנחש אותו.
דרישת הדיוק (\en{REQ-02}) גבוהה מן הדיוק שדווח באב-טיפוס מבוסס מצלמה, \pct{72} בספירה \cite{wickramasinghe2020}, ולכן בחירת החיישן נבדקת במעבדה לפני ההתקנה.
דרישת הפרטיות (\en{REQ-11}) קובעת את הארכיטקטורה: התמונה, אם יש בכלל מצלמה, אינה יוצאת מן החיישן (פרק~\ref{chap:architecture}).
